% ===========================================================================
%  CAPITULO 3 — Produto e Pedagogia (P01–P10)
% ===========================================================================
\chapter{Produto e Pedagogia}
\label{cap:produto}

% ---------------------------------------------------------------------------
\section{P01 · O que o professor vê antes de aplicar a missão}
\label{sec:p01}

\begin{resposta}{P01}{Ficha e guia da missão}{\implementado}
\begin{campos}
  \item[Hoje] Duas telas por missão, ambas no painel do professor, seção
    \emph{Entre Redes}, escolhendo ano e unidade:
    \begin{itemize}
      \item \textbf{Ficha da missão} (``Conhecer a missão''), em quatro abas:
        \emph{Como funciona} (os atos da missão), \emph{Antes de rodar}
        (preparo), \emph{O que ela não conta} (opcional) e \emph{Competência e
        formato}. Campos obrigatórios em toda missão publicada: título, resumo,
        unidade, competência, duração, formato, atos e preparo. Os códigos BNCC
        aparecem como etiquetas no cabeçalho (\adrr{ADR-100}, \adrr{ADR-182}).
      \item \textbf{Guia da aula}, em quatro tempos (\emph{Preparar},
        \emph{Conduzir}, \emph{Fechar}, \emph{Consultar}) e versão para
        impressão. Campos obrigatórios: objetivos; habilidades BNCC (código,
        texto, eixo e onde acontece na missão); preparo técnico e pedagógico e
        ``o que não dizer''; roteiro minuto a minuto com o comando de cada
        momento; debate, observação e mediação; síntese, registro,
        desdobramentos e \textbf{rubrica} (critério $\times$ em
        desenvolvimento / adequado / avançado); cuidados, adaptações e dúvidas
        frequentes. Avisos e produto final são opcionais. Um teste automático
        recusa publicar missão sem guia (\adrr{ADR-65}).
    \end{itemize}
    Os critérios por nível de cada habilidade aparecem também no painel de
    habilidades da atividade, nas missões que têm régua de avaliação.
  \item[Limitações] \textbf{Não há campo de bimestre}: a missão é identificada
    por ano/série e \emph{unidade} (1 a 4), e a equivalência unidade =
    bimestre não está formalizada. \textbf{Não há campo de correspondência
    com o livro}: aparece como texto livre em alguns guias (por exemplo,
    ``capítulo 1 do livro, pp. 2--9'' no guia da 2ª série do EM). Materiais
    necessários ficam dentro do preparo, sem campo próprio. Nada disso é
    editável pelo professor: é conteúdo fixo, versionado pela Unooria.
  \item[Evidência] \arq{apps/\B{}hub/\B{}lib/\B{}entre-redes/\B{}catalogo.ts} (ficha);
    \arq{apps/\B{}hub/\B{}lib/\B{}entre-redes/\B{}guias/\B{}} (um guia por missão) e o teste
    \arq{guias/index.test.ts}; suíte \arq{scripts/\B{}e2e/\B{}professor/\B{}08-guias-de-aula.mjs};
    Mapa de Produto, H04.3.2 e H04.4.1.
  \item[Lacuna] Produto: capturas das duas telas de uma missão real para anexo
    (podem ser tiradas da suíte e2e citada). Pedagogia/Editorial: formalizar
    unidade = bimestre e a referência ao livro como campo.
\end{campos}
\end{resposta}

% ---------------------------------------------------------------------------
\section{P02 · O que o professor informa para contextualizar}
\label{sec:p02}

\begin{resposta}{P02}{Formulário de personalização}{\variapormissao}
\begin{campos}
  \item[Hoje] Na ficha, \emph{Nova atividade} abre um assistente de quatro
    passos: \textbf{Turma \seta{} Conteúdo \seta{} Equipes \seta{} Revisar e
    criar} (\adrr{ADR-130}). O passo \emph{Conteúdo} é próprio de cada missão;
    as entradas são texto livre e escolha do componente curricular. Na maioria
    das missões há o atalho \emph{Usar banco padrão}, que carrega conteúdo
    curado pela Unooria; sem saldo de IA ou com o provedor indisponível, a
    tela usa o banco padrão. A Tabela~\ref{tab:p02} resume as entradas e o que
    a IA produz.
  \item[Fixo] Em todas as missões a IA \textbf{não altera} mecânica, etapas,
    gabaritos, rubricas e réguas de habilidade. Exemplos: em \emph{Quem Mandou
    Isso?} a resposta certa de cada mensagem não é editável; no \emph{Porta-Voz
    dos Dados} os números da série e o laudo nem passam pelo modelo; no
    \emph{Enigma da Capa} é o professor quem marca na imagem onde estão os
    objetos.
  \item[Limitações] \textbf{Não é possível anexar documentos} (PDF, plano de
    aula, apostila ou link) para orientar a geração: não existe campo de
    envio de arquivo. Os campos ``nome do arquivo'' que aparecem em duas
    missões são conteúdo fictício da simulação, não anexos. Duas missões não
    têm personalização por IA (\emph{Rastro do Boato} e a versão anterior da
    2ª série do EM).
  \item[Evidência] \arq{apps/\B{}hub/\B{}components/\B{}atividades/\B{}assistente-de-atividade.tsx};
    entradas reais das suítes e2e: ``Boato da merenda cancelada na escola''
    (\arq{estudante/\B{}10-quem-mandou-isso.mjs}), ``Água limpa e saúde no
    bairro'' com quatro tópicos (\arq{estudante/\B{}12-jornal-da-escola.mjs}),
    ``Revolução Industrial na Inglaterra'' (\arq{professor/\B{}02-atividade-do-zero.mjs}).
  \item[Lacuna] O resultado gerado não fica gravado nas suítes (é gerado a
    cada execução). Produto: anexar um par entrada/saída capturado em
    homologação.
\end{campos}
\end{resposta}

\begin{center}
\footnotesize\renewcommand{\arraystretch}{1.3}
\begin{tabularx}{\linewidth}{@{}L{34mm}L{58mm}X@{}}
\caption{Entradas do professor e o que a IA gera, por missão}\label{tab:p02}\\
\toprule
\textbf{Missão} & \textbf{O professor informa} & \textbf{A IA gera}\\
\midrule
\endfirsthead
\toprule
\textbf{Missão} & \textbf{O professor informa} & \textbf{A IA gera}\\
\midrule
\endhead
O Que a Foto Conta (1º EF) & tema fechado: pátio, horta, sala ou parquinho
  & legendas e comentários\\
Enigma da Capa (2º EF) & componente, tema, 3 a 6 palavras-chave
  & a cena e a imagem panorâmica\\
Nina e a Notícia (3º EF) & componente, tema do evento, assunto
  & o texto das 3 notícias (estrutura fixa)\\
Jornal da Escola (4º EF) & componente, tema central, tópicos das 8 notícias
  & 8 notícias e imagens\\
Quem Mandou Isso? (5º EF) & componente, boato fictício, 3 a 6 fatos, contexto
  local (opcional) & 8 mensagens\\
Júri Estadual (6º EF) & componente, título da mostra, 3 a 6 assuntos
  & 40 estandes e imagens\\
A Seleção (7º EF) & evento central, detalhe da cidade (opcional)
  & 12 posts e as falas do podcast\\
Unoovista (8º EF) & título, disciplina, tema, o que já foi trabalhado
  (opcional), proporção & perguntas-chave com nota de equilíbrio; posts e
  imagens\\
Manifesto da IA (8º EF) & dilema da escola (opcional) & textos de apoio\\
Dossiê Invisível (9º EF) & componente, tema, conceitos, contexto (opcional)
  & documentos do dossiê (quantidade fixa)\\
Assistente Invisível (1ª EM) & disciplina, tema, conceitos, foco regional
  (opcional) & as três conversas do assistente\\
Efeito Amplificador (2ª EM) & componente, tema da campanha, 3 a 6 tópicos,
  contexto & conteúdo da campanha\\
Porta-Voz dos Dados (3ª EM) & componente, tema gerador, 3 a 6 objetos de
  conhecimento, bairro (opcional) & textos (números e laudo fixos)\\
Rastro do Boato (6º EF) & --- & sem IA; conteúdo fixo\\
\bottomrule
\end{tabularx}
\end{center}

% ---------------------------------------------------------------------------
\section{P03 · Salvar, reaproveitar e compartilhar}
\label{sec:p03}

\begin{resposta}{P03}{A missão contextualizada como versão reutilizável}{\parcial}
\begin{campos}
  \item[Hoje] A atividade criada fica \textbf{salva} e pertence a
    \textbf{uma turma}. Na mesma turma ela é reaproveitável: a lista
    \emph{Suas atividades} mostra as atividades em andamento, prontas e em
    preparo, e \emph{Usar esta atividade agora} a põe no ar de novo, abrindo
    uma rodada nova e preservando a anterior. Só uma atividade fica no ar por
    turma; a que sai volta a ``pronta'' com tudo guardado. Também é possível
    reaproveitar as equipes da atividade anterior.
  \item[Quem faz o quê] Vê e conduz: o professor que leciona na turma (outra
    turma responde como inexistente). Pela interface, só o professor edita;
    a coordenação lê. Exclusão só existe numa missão (Unoovista), com
    confirmação pela digitação do título. Não há fluxo de aprovação por
    terceiros. O conteúdo aprovado fica gravado na própria atividade; se a
    Unooria atualizar a missão-base, a atividade já criada mantém o conteúdo
    aprovado.
  \item[Limitações] \textbf{Duplicar para outra turma, outro período ou outro
    professor não existe atualmente} (exceção: no Unoovista a atividade pode
    ser limpa e usada com outra turma). \textbf{Compartilhar entre
    professores, escola, rede ou biblioteca comum não existe atualmente.} No
    exemplo do questionário, a professora refaz a personalização para o
    8º ano B; o banco padrão da Unooria é o único conteúdo comum.
  \item[Planejado] Nada registrado no Mapa de Produto nem no \emph{backlog}.
  \item[Evidência] Mapa de Produto, H06.13.1 e H06.13.3; \adrr{ADR-130},
    \adrr{ADR-146}, \adrr{ADR-155}.
\end{campos}
\end{resposta}

% ---------------------------------------------------------------------------
\section{P04 · Curadoria do conteúdo gerado}
\label{sec:p04}

\begin{resposta}{P04}{Revisão antes de a turma ver}{\variapormissao}
\begin{campos}
  \item[Hoje] Regra comum: \textbf{nada gerado chega à turma sem aprovação do
    professor}; sem ela, a plataforma recusa pôr a atividade no ar. O
    professor vê cada item gerado. Os controles variam:
    \begin{itemize}
      \item \textbf{Editar e regenerar item a item}, com aprovação final do
        pacote inteiro: Assistente Invisível, Porta-Voz dos Dados e Dossiê
        Invisível (que marca ``modificado pelo docente''); na Nina só a
        manchete é editável.
      \item \textbf{Duas fases}, textos e depois imagens (\adrr{ADR-146},
        \adrr{ADR-147}): no Jornal da Escola cada imagem é aprovada
        individualmente e ``gerar outra imagem'' a reabre; na Seleção a troca é
        por post e a aprovação das imagens, em lote.
      \item \textbf{Aprovação em lote}: Júri Estadual e imagens da Nina; depois
        de aprovada, a imagem não muda.
      \item \textbf{Unoovista}: filtrar, regerar texto ou imagem, excluir post;
        ativar é recusado se faltar conteúdo.
      \item As falas do podcast são revisadas obrigatoriamente na criação
        (\adrr{ADR-150}, \adrr{ADR-204}).
    \end{itemize}
    Rejeitar um item é pedir nova geração manual daquele item; não há nova
    geração automática. \textbf{Registro}: a maioria das missões grava quem
    aprovou e quando, e o conteúdo aprovado fica salvo na atividade.
    \textbf{Verificação adicional}: regras automáticas conferem estrutura e
    critérios pedagógicos da saída antes de mostrá-la; na Seleção, uma
    segunda leitura por IA confere se o post editado manteve o enquadramento.
  \item[Limitações] Três missões não guardam o autor da aprovação (Nina guarda
    só a data; Manifesto, um indicador sim/não; Unoovista, o dono e o
    estado). Edições não são versionadas. Não há moderação dedicada nem
    configuração própria dos filtros do provedor (T05). O Mapa de Produto
    registra que as imagens geradas precisam de revisão de conteúdo antes da
    primeira rede (H06.12.7).
  \item[Evidência] \arq{apps/\B{}hub/\B{}app/\B{}(plataforma)/\B{}professor/\B{}entre-redes/\B{}AGENTS.md};
    \cd{PU-DOC-001}, Cap. 11; suítes e2e de estudante 10 a 19 e
    \arq{professor/\B{}02-atividade-do-zero.mjs}.
\end{campos}
\end{resposta}

% ---------------------------------------------------------------------------
\section{P05 · Comandos do professor durante a aula}
\label{sec:p05}

\begin{resposta}{P05}{Condução em sala}{\variapormissao}
\begin{campos}
  \item[Hoje] Em toda missão há três comandos de base: \textbf{liberar} a
    missão na turma (sem isso, o estudante vê o cartão com cadeado, e
    digitar o endereço também não abre: \adrr{ADR-69}), \textbf{pôr a
    atividade no ar} e \textbf{encerrar} (fecha a entrada sem expulsar quem
    está jogando, e converte as marcas pendentes em nível). Depois, cada
    missão tem as etapas próprias, conduzidas pelo professor
    (Tabela~\ref{tab:p05}). O estudante vê a tela de espera, a tela
    congelada (pausa) ou a etapa nova; quem está atrasado continua no próprio
    caminho. O painel da aula mostra ``ainda não entrou'', ``concluíram
    $x$/$y$'', a fila por passo e o tempo parado; com equipes, conta por
    computador.
  \item[Limitações] A atualização é por consulta periódica, não em tempo
    real: o estudante relê a etapa a cada 3 a 5 segundos e o painel do
    professor a cada 2 a 5 segundos (\adrr{ADR-93}); todos precisam estar
    conectados. \textbf{As etapas só avançam}: voltar ou pular é recusado, e
    repetir o comando não muda nada. Programar abertura e fechamento ainda não
    existe (H04.2.3, a fazer).
  \item[Evidência] Uma suíte e2e por missão em \arq{scripts/e2e/estudante/};
    \arq{scripts/\B{}e2e/\B{}professor/\B{}06-resultados-das-missoes.mjs}; Mapa de
    Produto, H04.2.1 e H04.2.2.
\end{campos}
\end{resposta}

\begin{center}
\footnotesize\renewcommand{\arraystretch}{1.3}
\begin{tabularx}{\linewidth}{@{}L{42mm}X@{}}
\caption{Comandos de condução, por missão}\label{tab:p05}\\
\toprule
\textbf{Missão} & \textbf{Etapas e comandos do professor}\\
\midrule
\endfirsthead
\toprule
\textbf{Missão} & \textbf{Etapas e comandos do professor}\\
\midrule
\endhead
O Que a Foto Conta & configurar, aprovar o pacote, destravar etapa, reiniciar\\
Enigma da Capa & liberar a virada\\
Nina e a Notícia & aposta \seta{} filtro \seta{} inspeção; pausar e retomar a turma\\
Jornal da Escola & exploração \seta{} resultado \seta{} revisão\\
Quem Mandou Isso? & revelação, síntese (opcional), reflexão; projetar o termômetro da turma\\
Júri Estadual & exploração \seta{} galeria (liberar auditoria)\\
Rastro do Boato & liberar a aula 2; pausar a simulação\\
A Seleção & exploração \seta{} mural \seta{} fichas \seta{} explicação\\
Manifesto da IA & pausar para debate; retomar\\
Dossiê Invisível & exploração \seta{} revelação \seta{} revoto \seta{} explicação (sem pausa)\\
Assistente Invisível & testes \seta{} simulação \seta{} explicar \seta{} publicar, com duas pausas
  de debate opcionais; a simulação exige todas as apostas feitas\\
Efeito Amplificador & exploração \seta{} pausa para debate \seta{} explicar\\
Porta-Voz dos Dados & exploração \seta{} pausa \seta{} liberação da errata\\
\bottomrule
\end{tabularx}
\end{center}

% ---------------------------------------------------------------------------
\section{P06 · Nova aplicação e novas tentativas}
\label{sec:p06}

\begin{resposta}{P06}{Missão-base, atividade e rodada}{\implementado}
\begin{campos}
  \item[Hoje] Os três objetos existem: a \textbf{missão-base} (o esqueleto
    publicado pela Unooria), a \textbf{atividade} (a aplicação
    contextualizada, de uma turma) e a \textbf{rodada} (cada vez que a
    atividade vai ao ar). (1) O professor pode criar \textbf{várias atividades
    da mesma missão para a mesma turma}; só uma fica no ar por vez.
    (2) Pôr no ar de novo uma atividade que já rodou abre a \textbf{rodada
    seguinte}, e o estudante refaz a missão; a rodada anterior fica guardada
    e toda rodada mede tudo de novo (\adrr{ADR-145}). \emph{Zerar a turma}
    (\adrr{ADR-155}) abre nova rodada e apaga as marcas e os níveis daquela
    atividade, inclusive ajustes do professor, para recomeçar limpo (útil
    após erro); o histórico de eventos permanece. (3) \textbf{O que vale}: na
    mesma rodada, o melhor registro; na atividade, a melhor rodada por
    habilidade; no painel de habilidades, o melhor resultado entre todas as
    atividades e rodadas da missão (\adrr{ADR-144}). Ajuste do professor
    prevalece sobre registro próprio, que prevalece sobre o da equipe.
    (4) \textbf{Duplicidade}: o banco aceita um único resultado por escola,
    missão, atividade, rodada, estudante e habilidade; liberar e pôr no ar são
    idempotentes; atividade e rodada nunca vêm do navegador, são definidas no
    servidor.
  \item[Limitações] Exceções: no Unoovista vale o primeiro envio; em
    múltipla escolha vale a primeira resposta, salvo régua que diga o
    contrário. Não há regra específica para ano letivo seguinte: a atividade
    é da turma, e a turma nova é outra turma (P03).
  \item[Evidência] Mapa de Produto, H14.2.2 e H14.2.4; testes
    \arq{apps/\B{}api/\B{}tests/\B{}evidencias-consolidado-postgres.spec.ts}; suítes
    \arq{professor/\B{}13-painel-de-habilidades.mjs} e
    \arq{professor/\B{}11-perfil-do-aluno.mjs}.
\end{campos}
\end{resposta}

% ---------------------------------------------------------------------------
\section{P07 · Evidência quando o estudante conversa com a IA}
\label{sec:p07}

\begin{resposta}{P07}{Conversa livre e avaliação}{\implementado}
\begin{campos}
  \item[Hoje] Três missões têm texto livre do estudante dirigido à IA: o
    \emph{Efeito Amplificador} (2ª EM), o \emph{Dossiê Invisível} (9º EF) e a
    versão anterior da 2ª série (fora do catálogo, mantida para turmas que já
    a usavam). O Assistente Invisível (1ª EM) \textbf{não} tem conversa
    livre: o estudante lê conversas preparadas. (1) \textbf{A avaliação usa
    eventos estruturados, nunca o conteúdo da conversa}: o que pontua são
    ações (publicar, avisar, proteger, escolher um caminho). Nas missões
    atuais, a IA só traduz a mensagem numa intenção de lista fechada, que
    decide o caminho da história; ``a IA nunca decide o que a régua lê''
    (\adrr{ADR-192}). (2) Os critérios estão nas \textbf{réguas de
    avaliação}, uma por missão, escritas a partir da régua pedagógica, com
    códigos BNCC conferidos automaticamente contra o currículo e versão
    gravada em cada marca. (3) Há LLM classificando em dois pontos: a
    intenção da mensagem (temperatura 0, saída em formato fechado e validado,
    confiança mínima de 0,7, árvore de regras quando a IA falha) e, na
    \emph{Seleção}, a legenda escrita pelo estudante, lida em três
    indicadores; o nível é decidido por regra no servidor e, se a IA falhar,
    fica ``não avaliado'' para o professor ajustar. A consistência é medida
    por conjunto de referência com a IA real: 88/89 e 89/89 acertos, zero
    crédito indevido. (4) O texto do estudante não vai a painel, telemetria
    nem registro técnico; na versão anterior da 2ª série o professor lê a
    transcrição de cada estudante, e o gestor só vê o agregado.
  \item[Exemplo] Do conjunto de referência: a mensagem ``essa frase da empatia
    é sua mesmo?'' é classificada como a intenção ``auditar autoria'', que
    leva ao caminho em que o estudante investiga a origem do texto; a
    evidência vem desse caminho e das ações seguintes, não da frase.
  \item[Limitações] O classificador não é determinístico: um erro custa ao
    estudante reformular a mensagem, nunca um crédito indevido.
  \item[Evidência] \arq{apps/\B{}api/\B{}src/\B{}efeito-amplificador/\B{}AGENTS.md};
    \arq{apps/\B{}api/\B{}scripts/\B{}medir-classificador-exp02em1b.ts};
    \arq{apps/\B{}api/\B{}src/\B{}evidencias/\B{}fichas/\B{}}; \arq{tests/\B{}evidencias-fichas.spec.ts};
    \adrr{ADR-172}, \adrr{ADR-191}--\adrr{ADR-193}.
\end{campos}
\end{resposta}

% ---------------------------------------------------------------------------
\section{P08 · Qualidade das recomendações pedagógicas}
\label{sec:p08}

\begin{resposta}{P08}{Recomendações geradas por IA}{\naoexiste}
\begin{campos}
  \item[Hoje] \textbf{Recomendações de intervenção geradas por IA a partir das
    dificuldades da turma não existem atualmente.} O que existe: o
    \textbf{painel de habilidades} mostra, por habilidade, quantos estudantes
    atingiram, estão em processo, não atingiram ou não foram avaliados, e
    ordena por regra fixa (sem IA) para o professor ver onde reforçar; numa
    missão a IA \textbf{sugere um nível} para a reflexão do estudante,
    rotulado ``sugestão da IA'' até o professor confirmar (\adrr{ADR-167});
    o \textbf{agente de IA do professor} responde sobre BNCC e sobre o acervo
    de livros, citando código e página, mas não acessa resultados da turma
    nem evidências.
  \item[Prompts] Ficam no repositório, só no servidor, versionados pelo
    histórico do código; apenas uma missão grava a versão do prompt junto do
    resultado. A autoria pedagógica dos prompts não está registrada nos
    arquivos.
  \item[Observado $\times$ sugerido] Já há separação explícita: sugestão da IA
    é rotulada e só vale após confirmação; ``não avaliado'' nunca conta como
    ``não atingiu''; o agente só cita habilidade que a busca devolveu e avisa
    quando o código não existe.
  \item[Limitações] O agente do professor está parcial e ainda não está no ar
    nos ambientes da rede.
  \item[Evidência] \arq{apps/\B{}hub/\B{}lib/\B{}painel-de-habilidades.ts};
    \arq{apps/\B{}agente-professor/\B{}agente\_professor/\B{}instrucoes.py} e o conjunto de
    avaliação \arq{avaliacao/\B{}agente\_professor.evalset.json} (16 casos);
    Mapa de Produto, F08.4 e H14.4.5.
  \item[Lacuna] É o compromisso V1-03/V1-04 (Capítulo~\ref{cap:v1}).
\end{campos}
\end{resposta}

% ---------------------------------------------------------------------------
\section{P09 · Acessibilidade e inclusão}
\label{sec:p09}

\begin{resposta}{P09}{Recursos por tipo de missão}{\parcial}
\begin{campos}
  \item[Hoje] Alvo declarado: \textbf{WCAG 2.2 AA} (\adrr{ADR-162}), com
    verificação automática (axe-core) nas regras A e AA das versões 2.0, 2.1
    e 2.2, medindo a página e cada missão embutida. Primeira medição: 65 telas
    do hub e a abertura de 11 missões; as violações encontradas foram
    corrigidas no design system (contraste de texto, foco visível, alvo de
    toque, relações de acessibilidade das abas), e a última medição
    registrada é de \textbf{zero violação} nessas telas. Recursos presentes:
    \textbf{texto alternativo} nas imagens geradas por IA, derivado da
    descrição que gerou a imagem, nunca vazio (\adrr{ADR-163}), e nas fotos
    fixas; \textbf{transcrição integral sempre visível} nos áudios de
    podcast, pensada também para sala sem fone; anúncio de mudança para
    leitor de tela em várias missões; respeito à preferência de movimento
    reduzido em várias missões. \textbf{1º ano, não alfabetizados}: instrução
    em caixa alta com até 8 palavras e redundância tripla (texto, pictograma e
    locução), botão grande de repetir o áudio, uma locução por vez.
  \item[Limitações] Só houve verificação automática: \textbf{falta a passada
    manual com teclado e leitor de tela}, que é a que autoriza declarar AA
    (risco R-07). A Vila do Lume (jogo 3D) está fora da medição. As
    \textbf{locuções ainda não foram gravadas} (a missão do 1º ano cai para
    texto e ícone; Manifesto e Rastro do Boato aguardam áudio). \textbf{Não
    há controle para ligar e desligar áudio}: ouvir é sempre opcional (o
    áudio só toca por comando) e o texto sempre o substitui. Não há legenda de
    vídeo nem ampliação de fonte própria (vale o zoom do navegador). A
    preferência de movimento reduzido não está em todas as missões.
  \item[Planejado] Passada manual de teclado e leitor de tela por papel;
    tornar a verificação bloqueante no catálogo de componentes; gravação das
    locuções (responsável: Conteúdo).
  \item[Evidência] \arq{scripts/\B{}e2e/\B{}lib/\B{}acessibilidade.mjs};
    \arq{scripts/\B{}e2e/\B{}ferramentas/\B{}medir-acessibilidade.mjs};
    \cd{PU-DOC-001}, Cap. 16 (estado ``parcial'') e Cap. 19 (R-07).
  \item[Lacuna] O relatório da última medição não é guardado no repositório;
    gerar e anexar um novo antes do TR.
\end{campos}
\end{resposta}

\begin{paraotr}
A solução deve atender à WCAG 2.2 nível AA (ou ABNT NBR 17225) nas telas de
estudante, professor e gestor, com alternativa textual a todo áudio e imagem.
Aceite: relatório de verificação automática sem violações e relatório de
avaliação manual com teclado e leitor de tela, com plano de correção datado
para as barreiras encontradas.
\end{paraotr}

% ---------------------------------------------------------------------------
\section{P10 · Resultados da PoC com o Instituto Paula Souza}
\label{sec:p10}

\begin{resposta}{P10}{PoC}{\foradorepo}
\begin{campos}
  \item[Hoje] O repositório não contém relatório, apresentação, instrumento de
    pesquisa nem dados da PoC; não há menção ao Instituto Paula Souza no
    código nem na documentação. Também não há instrumentação de pesquisa
    (questionário de satisfação, medida de aprendizagem pré e pós) na
    plataforma.
  \item[O que a plataforma oferece] Para uma futura avaliação, os dados que a
    plataforma já registra podem servir de medida objetiva de conclusão e de
    desempenho por habilidade: progresso e conclusão por estudante e missão,
    evidências por habilidade e telemetria das missões (T10).
  \item[Lacuna] Time de Produto/Pedagogia: anexar período, objetivos,
    missões testadas, número e perfil de participantes, metodologia,
    instrumentos, resultados, limitações e mudanças de produto. Sem esse
    material, o TR não deve declarar ganho de aprendizagem, engajamento ou
    aceitação.
\end{campos}
\end{resposta}
